\documentclass[a4paper]{article}
\usepackage{ctex}
\usepackage[affil-it]{authblk}
\usepackage[backend=bibtex,style=numeric]{biblatex}
\usepackage{amsmath,amsthm,amssymb,amsfonts}
\usepackage{graphicx}
\usepackage{bm}

\usepackage{geometry}
\geometry{margin=1.5cm, vmargin={0pt,1cm}}
\setlength{\topmargin}{-1cm}
\setlength{\paperheight}{29.7cm}
\setlength{\textheight}{25.3cm}

\addbibresource{citation.bib}

\begin{document}
% =================================================
\title{微分方程数值解 2026 春夏作业三}

\author{曾申昊 3220100701
  \thanks{邮箱: \texttt{3097714673@qq.com}
                            \texttt{或} \texttt{3220100701@zju.edu.cn}}}
\affil{浙江大学}

\date{更新时间: \today}

\maketitle

% ============================================
\section*{Exercise 12.14}

\par 初值问题的积分形式为：
\begin{gather}
    \mathbf{v}(t) = \mathbf{v}_0 
    + \int_a^t \mathbf{f}( , \mathbf{v}(s)) \mathrm{d}s\\
    \mathbf{w}(t) = \mathbf{w}_0
    + \int_a^t \mathbf{f}(s, \mathbf{w}(s)) \mathrm{d}s
\end{gather}
\par 两式相减并两边取范数，得到：
\begin{equation}
    \begin{aligned}
        \|\mathbf{v}(t) - \mathbf{w}(t)\| 
        &\leq \|\mathbf{v}_0 - \mathbf{w}_0\| 
        + \int_a^t \|\mathbf{f}(s, \mathbf{v}(s)) - \mathbf{f}(s, \mathbf{w}(s))\| \mathrm{d}s\\
        &\leq \|\mathbf{v}_0 - \mathbf{w}_0\| 
        + L\int_a^t \|\mathbf{v}(s) - \mathbf{w}(s)\| \mathrm{d}s
    \end{aligned}
\end{equation}
\par 由 Gronwall 不等式，得到：
\begin{equation}
    \|\mathbf{v}(t) - \mathbf{w}(t)\|
    \leq \|\mathbf{v}_0 - \mathbf{w}_0\| e^{L(t-a)}
\end{equation}

\section*{Exercise 12.64}

\par 梯形法的迭代公式为：
\begin{equation}
    \mathbf{U}^{n+1}
    -\mathbf{U}^{n}
    =k\left(
        \frac{1}{2} \mathbf{f}\left(\mathbf{U}^{n+1}, t_{n+1}\right)
        +\frac{1}{2} \mathbf{f}\left(\mathbf{U}^{n}, t_{n}\right)
    \right)
\end{equation}
\par 进行 Taylor 展开，得到：
\begin{equation}
    \begin{aligned}
        \mathcal{L}\mathbf{u}(t_n)
        =&\ \mathbf{u}(t_{n+1})
        -\mathbf{u}(t_{n})
        -k\left(
            \frac{1}{2} \mathbf{u}'(t_{n+1})
            +\frac{1}{2} \mathbf{u}'(t_{n})
        \right)\\
        \approx&\left(
            u(t_n) + k u'(t_n) + \frac{k^2}{2} u''(t_n) + \frac{k^3}{6} u'''(t_n) + \frac{k^4}{24} u^{(4)}(t_n)
        \right)
        -u(t_n)\\
        &-k\left[
            \frac{1}{2} \left(
                u'(t_n) + k u''(t_n) + \frac{k^2}{2} u'''(t_n) + \frac{k^3}{6} u^{(4)}(t_n)
            \right)
            +\frac{1}{2} u'(t_{n})
        \right]\\
        = &\  \frac{k^3}{12} u'''(t_n) + \frac{k^4}{24} u^{(4)}(t_n)
    \end{aligned}
\end{equation}
\par 因此 $C_0=C_1=C_2=0$，$C_3=\dfrac{1}{12}$，$C_4=\dfrac{1}{24}$。
\par 中点法的迭代公式为：
\begin{equation}
    \mathbf{U}^{n+2}
    -\mathbf{U}^{n}
    =k\left(
        2\mathbf{f}\left(\mathbf{U}^{n+1}, t_{n+1}\right)
    \right)
\end{equation}
\par 进行 Taylor 展开，得到：
\begin{equation}
    \begin{aligned}
        \mathcal{L}\mathbf{u}(t_n)
        =&\ \mathbf{u}(t_{n+2})
        -\mathbf{u}(t_{n})
        -k\left(
            2 \mathbf{u}'(t_{n+1})
        \right)\\
        \approx&\left(
            u(t_n) + 2k u'(t_n) + \frac{(2k)^2}{2} u''(t_n) 
            + \frac{(2k)^3}{6} u'''(t_n) + \frac{(2k)^4}{24} u^{(4)}(t_n)
        \right)
        -u(t_n)\\
        &-k\left[
            2 \left(
                u'(t_n) + k u''(t_n) + \frac{k^2}{2} u'''(t_n) + \frac{k^3}{6} u^{(4)}(t_n)
            \right)
        \right]\\
        = &\  \frac{k^3}{3} u'''(t_n) + \frac{k^4}{3} u^{(4)}(t_n)
    \end{aligned}
\end{equation}
\par 因此 $C_0=C_1=C_2=0$，$C_3=\dfrac{1}{3}$，$C_4=\dfrac{1}{3}$。

\section*{Exercise 12.66}

\par $\left\|\mathcal{L}\mathbf{u}(t_n)\right\|=O(k^3)$ 等价于：
\begin{gather}
    \sum_{j=0}^{s}\alpha_j = 0\\
    \sum_{j=0}^{s} j \alpha_j = \sum_{j=0}^{s} \beta_j\\
    \frac{1}{2}\sum_{j=0}^{s} j^2 \alpha_j = \sum_{j=0}^{s} j \beta_j
\end{gather}
\par 用特征多项式表示为：
\begin{gather}
    \rho''(\zeta) = \sum_{j=0}^{s} j(j-1)\alpha_j \zeta^{j-2}
    =\sum_{j=0}^{s} j^2\alpha_j \zeta^{j-2}-\sum_{j=0}^{s} j \alpha_j \zeta^{j-2}\\
    \sigma''(\zeta) = \sum_{j=0}^{s} j(j-1)\beta_j \zeta^{j-2}
    =\sum_{j=0}^{s} j^2\beta_j \zeta^{j-2}-\sum_{j=0}^{s} j \beta_j \zeta^{j-2}
\end{gather}
\par 等价于：
\begin{gather}
    \rho(1) = 0\\
    \rho'(1) = \sigma(1)\\
    \frac{1}{2}\left(
        \rho''(1)
        +\rho'(1)
    \right) = \sigma'(1)
\end{gather}

\section*{Exercise 12.67}

\par 对于 Adams-Bashforth 法有 $\alpha_s=1, \alpha_{s-1}=-1, 
\alpha_1=\cdots=\alpha_{s-2}=0$ 以及 
$\beta_s=0$。下面以 $s=3, p=3$ 为例进行待定系数法计算：
\begin{equation}
    \begin{bmatrix}
        1 & 0 & 0 & 0 & 0 & 0 & 0 & 0\\
        0 & 1 & 0 & 0 & 0 & 0 & 0 & 0\\
        0 & 0 & 1 & 0 & 0 & 0 & 0 & 0\\
        0 & 0 & 0 & 1 & 0 & 0 & 0 & 0\\
        0 & 1 & 2 & 3 & -1 & -1 & -1 & -1\\
        0 & \frac{1^2}{2} & \frac{2^2}{2} & \frac{3^2}{2} & 0 & -1 & -2 & -3\\
        0 & \frac{1^3}{6} & \frac{2^3}{6} & \frac{3^3}{6} & 0 & -\frac{1^2}{2} & -\frac{2^2}{2} & \frac{3^2}{2}\\
        0 & 0 & 0 & 0 & 0 & 0 & 0 & 1\\
    \end{bmatrix}
    \begin{bmatrix}
        \alpha_0\\
        \alpha_1\\
        \alpha_2\\
        \alpha_3\\
        \beta_0\\
        \beta_1\\
        \beta_2\\
        \beta_3\\
    \end{bmatrix}
    =
    \begin{bmatrix}
        0\\
        0\\
        -1\\
        1\\
        0\\
        0\\
        0\\
        0\\
    \end{bmatrix}
\end{equation}
\par 计算可得 $\beta_3=0$，$\beta_2=\dfrac{23}{12}$，
$\beta_1=-\dfrac{16}{12}$，$\beta_0=\dfrac{5}{12}$。
\par 对于 Adams-Moulton 法有 $\alpha_s=1, \alpha_{s-1}=-1, 
\alpha_1=\cdots=\alpha_{s-2}=0$
以及 $\beta_s\neq 0$。下面以 $s=2, p=3$ 为例进行待定系数法计算：
\begin{equation}
    \begin{bmatrix}
        1 & 0 & 0 & 0 & 0 & 0 \\
        0 & 1 & 0 & 0 & 0 & 0 \\
        0 & 0 & 1 & 0 & 0 & 0 \\
        0 & 1 & 2 & -1 & -1 & -1 \\
        0 & \frac{1^2}{2} & \frac{2^2}{2}& 0 & -1 & -2\\
        0 & \frac{1^3}{6} & \frac{2^3}{6}& 0 & -\frac{1^2}{2} & -\frac{2^2}{2}\\
    \end{bmatrix}
    \begin{bmatrix}
        \alpha_0\\
        \alpha_1\\
        \alpha_2\\
        \beta_0\\
        \beta_1\\
        \beta_2\\
    \end{bmatrix}
    =
    \begin{bmatrix}
        0\\
        -1\\
        1\\
        0\\
        0\\
        0\\
    \end{bmatrix}
\end{equation}
\par 计算可得 $\beta_2=\dfrac{5}{12}$，
$\beta_1=\dfrac{8}{12}$，$\beta_0=-\dfrac{1}{12}$。
\par 对于 BDF 法有 $\alpha_3=1$ 以及 $\beta_s\neq 0, 
\beta_1=\cdots=\beta_{s-1}=0$。下面以 $s=3, p=3$ 为例进行待定系数法计算：
\begin{equation}
    \begin{bmatrix}
        1 & 1 & 1 & 1 & 0 & 0 & 0 & 0\\
        0 & 1 & 2 & 3 & -1 & -1 & -1 & -1\\
        0 & \frac{1^2}{2} & \frac{2^2}{2} & \frac{3^2}{2} & 0 & -1 & -2 & -3\\
        0 & 0 & 0 & 1 & 0 & 0 & 0 & 0\\
        0 & 0 & 0 & 0 & 1 & 0 & 0 & 0\\
        0 & 0 & 0 & 0 & 0 & 1 & 0 & 0\\
        0 & 0 & 0 & 0 & 0 & 0 & 1 & 0\\
        0 & \frac{1^3}{6} & \frac{2^3}{6} & \frac{3^3}{6} & 0 & -\frac{1^2}{2} & -\frac{2^2}{2} & \frac{3^2}{2}\\
    \end{bmatrix}
    \begin{bmatrix}
        \alpha_0\\
        \alpha_1\\
        \alpha_2\\
        \alpha_3\\
        \beta_0\\
        \beta_1\\
        \beta_2\\
        \beta_3\\
    \end{bmatrix}
    =
    \begin{bmatrix}
        0\\
        0\\
        0\\
        1\\
        0\\
        0\\
        0\\
        0\\
    \end{bmatrix}
\end{equation}
\par 计算可得 $\alpha_3=1$，$\alpha_2=-\dfrac{18}{11}$，
$\alpha_1=\dfrac{9}{11}$，$\alpha_0=-\dfrac{2}{11}$
以及 $\beta_3=\dfrac{6}{11}$。

\section*{Exercise 12.72}

\par 三阶精度的 BDF 法的特征多项式为：
\begin{gather}
    \rho(\zeta) = \sum_{j=0}^{s}\alpha_j \zeta^{j}
    =-\frac{2}{11} + \frac{9}{11} \zeta - \frac{18}{11} \zeta^2 + \zeta^3\\
    \sigma(\zeta) = \sum_{j=0}^{s}\beta_j \zeta^{j}
    =\frac{6}{11} \zeta^3
\end{gather}
\par 应用 \textbf{Theorem 12.70}可得
\begin{equation}
    \begin{aligned}
        \frac{\rho(z)}{\sigma(z)}
        =&\ \frac{-\frac{2}{11} + \frac{9}{11} z - \frac{18}{11} z^2 + z^3}{\frac{6}{11} z^3}\\
        =&-\frac{1}{3} \frac{1}{z^3} + \frac{3}{2} \frac{1}{z^2} - 3 \frac{1}{z} + \frac{11}{6}\\
        =&-\frac{1}{3} \frac{1}{(1+z-1)^3} + \frac{3}{2} \frac{1}{(1+z-1)^2} - 3 \frac{1}{1+z-1} + \frac{11}{6}\\
        =&-\frac{1}{3}\left(1-3(z-1)+6(z-1)^2-10(z-1)^3+15(z-1)^4+O((z-1))^5\right)\\
         &+\frac{3}{2}\left(1-2(z-1)+3(z-1)^2-4(z-1)^3+5(z-1)^4+O((z-1))^5\right)\\
         &-3\left(1-(z-1)+(z-1)^2-(z-1)^3+(z-1)^4+O((z-1))^5\right)+\frac{11}{6}\\
        =&\ \left(
            (z-1)-\frac{1}{2}(z-1)^2+\frac{1}{3}(z-1)^3-\frac{1}{2}(z-1)^4+O((z-1))^5
        \right)
    \end{aligned}
\end{equation}
\par 该 BDF 法的确是三阶精度的，
四阶误差常数为 $-\dfrac{1}{4}$。

\section*{Exercise 12.73}

\par “$\Rightarrow$” 的证明：
\par 由 \textbf{Definition 12.60} 可知一个 s-step 线性多步法
由 $p$ 阶数精度等价于 $\mathcal{L}\mathbf{u}(t_n)=O(k^{p+1})$。
对于任意 $\dfrac{d u}{dt}=q(t)$ ，$q(t)$ 为阶数小于 $p$ 的多项式，
都有 $\dfrac{d^{p+1} u}{dt^{p+1}}=q^{(p)}(t)=0$，
因此$\mathcal{L}\mathbf{u}(t_n)=0$，得到的解是精确的。

\par “$\Leftarrow$” 的证明：
\par 对于任意 $\dfrac{d u}{dt}=q(t)$ ，$q(t)$ 为阶数小于 $p$ 的多项式，
都有 $\dfrac{d^{p+1} u}{dt^{p+1}}=q^{(p)}(t)=0$。
对于任意函数 $u(t)$，用多项式逼近得到
\begin{equation}
    \mathcal{L}\mathbf{u}(t_n)
    =\mathcal{L}\mathbf{u}(t_n) - \mathcal{L}\mathbf{p}(t_n)\\
    =O(k^{p+1})
\end{equation}
\par 其中 $\mathbf{p}(t)$ 是用于逼近 $u(t)$ 的多项式，阶数小于 $p$。
因此该 s-step 线性多步法有 $p$ 阶数精度。

% ===============================================

\printbibliography

\end{document}